\documentclass[10pt,a4paper]{amsart}
\usepackage[margin=19mm]{geometry}
\usepackage{amsmath,amssymb,mathtools}
\usepackage[scr=rsfs,cal=euler]{mathalfa}
\usepackage{xfrac}
\usepackage[unicode,hidelinks,bookmarks=false]{hyperref}
\newcommand{\ie}{i.e.\ }
\newcommand{\cf}{cf.\ }
\newcommand{\resp}{respectively\ }
\newcommand{\Subsection}{\subsection}
\providecommand{\nobreakdash}{\nobreak\hbox{-}}
\setlength{\emergencystretch}{2em}
\multlinegap0pt
\usepackage{bm}
\usepackage{bbm}
\usepackage{dsfont}
\usepackage{xspace}
\usepackage{cancel}
\usepackage{mathtools}
\usepackage{amsthm}
\makeatletter
\@ifundefined{theorem}{\newtheorem{theorem}{Theorem}}{}
\@ifundefined{lemma}{\newtheorem{lemma}[theorem]{Lemma}}{}
\makeatother
\title[Constructing Parseval Fusion Frames via Operators]{Constructing Parseval Fusion Frames via Operators}
\author{Ehsan Ameli, Ali Akbar Arefijamaal, Fahimeh Arabyani Neyshaburi}
\date{Source: arXiv:2509.25108v1 (CC BY 4.0)}
\begin{document}
\maketitle

\noindent\emph{Source and licence.} Constructing Parseval Fusion Frames via Operators. Version arXiv:2509.25108v1 (CC BY 4.0), distributed under the Creative Commons Attribution 4.0 International licence, as recorded in the arXiv licence metadata for this exact version and shown on the abstract page https://arxiv.org/abs/2509.25108v1. Full rights notice: licenses/arxiv-2509.25108-CC-BY-4.0.txt.\par\smallskip

\noindent\textit{Editorial scope.} Counted unit: the Theorem whose source label is \texttt{in case Riesz} in the article (source0.tex lines 336--344, source label \texttt{in case Riesz}), delivered together with the article's own complete original proof of it (source0.tex lines 346--401, one \texttt{proof} environment of 56 lines). No mathematics is written, repaired or re-derived: statement and proof are the article's own text line for line. Required context retained uncounted: extend Gav (lines 218--220); r.e (lines 240--248). Every definition, display and label of those lines is retained unchanged. Omitted from this package and not redistributed: 1--217, 221--239, 249--335, 345--345, 402--1107. Nothing inside a retained excerpt is omitted or altered: every retained body is delivered exactly as the article's lines are written, so no omission receipt of any kind accompanies this package. Citations: the retained excerpts carry no citation command at all, so no bibliography entry of the article is transcribed and no reference is redistributed. Reference adaptation: none was needed and none was made; every reference inside the delivered unit resolves inside it, so no unresolved reference or citation is produced. Printed slips kept literal: no printed word, symbol, hypothesis or number of the counted statement or of its proof was altered. Layout and class: the article's own class is \texttt{sn-jnl} with packages including graphicx, multirow, amsmath, amssymb, amsfonts, amsthm, mathrsfs, appendix, xcolor, textcomp, manyfoot, booktabs, algorithm, algorithmicx; the wrapper is the lane's pinned \texttt{amsart} editorial wrapper at 10pt on A4, which restates the theorem-like environments the retained text needs and adds the article's own macro definitions that the retained text needs (none), with extra wrapper packages (bm, bbm, dsfont, xspace, cancel, mathtools). The delivered numbering is the unit's own. Assets: no retained span contains a figure environment, a graphics-inclusion command, a PDF-page inclusion, a table or a TikZ picture; no image file is copied into this package, the package declares no asset dependency and no image file is redistributed with it. Licence: version arXiv:2509.25108v1 of this article is distributed under the Creative Commons Attribution 4.0 International licence, as recorded in the arXiv licence metadata for this exact version and shown on the abstract page https://arxiv.org/abs/2509.25108v1. A full rights notice is delivered at licenses/arxiv-2509.25108-CC-BY-4.0.txt. Characters: every retained excerpt is pure ASCII, so the delivered engine, XeLaTeX with its default Latin Modern text font, reports no missing character.
\begin{lemma}\label{extend Gav}
Let $ W\subset \mathcal{H}$ be a closed subspace and $U,T \in B(\mathcal{H})$ be invertible operators such that $\pi_{UW}UT=U\pi_{W}.$ Then $U^{*}UW\subseteq \left( T^{*}\right) ^{-1}W.$
\end{lemma}

\begin{lemma}\label{r.e}
Let $\mathcal{W}=\{(W_{i},\omega_{i})\}_{i \in I} $ be a fusion frame for $\mathcal {H}$ and $U \in B(\mathcal{H})$ an invertible operator. Then the following are equivalent:
\begin{itemize}
\item[$(i)$] $\mathcal{W}$ is a fusion Riesz basis.
\item[$(ii)$] $ \left( U^*\right) ^{-1}S_{\mathcal{W}}^{-1}W_{i}\perp UW_{j}$ for all $i,j \in I, i\neq j. $
\item[$(iii)$] $ \omega_{i}^{2}\pi_{W_{i}}S_{W}^{-1}\pi_{W_{j}}=\delta_{i,j}\pi_{W_{j}}~ for~all~ i , j \in I $.
\item[$(iv)$] $\omega_{i}^{2}\pi_{UW_{i}}\left( U^*\right) ^{-1}S_{\mathcal{W}}^{-1}\pi_{W_{j}}=\delta_{i,j}\pi_{UW_{j}}\left( U^*\right) ^{-1}$ for all $i,j \in I.$
\end{itemize}
\end{lemma}

\begin{theorem}\label{in case Riesz}
Let $\mathcal{W}=\{(W_{i},\omega_{i})\}_{i \in I} $ be a fusion Riesz basis for $\mathcal {H}$. Then the following are equivalent:
\begin{itemize}
\item[$(i)$] $\mathcal{W}$ is $\left( U,\omega^{-1}\right) $-scalable.
\item[$(ii)$] $\mathcal{W}$ is $\left( \left( S_{\mathcal{W}}U^*\right) ^{-1},\omega^{-1}\right) $-scalable.
\item[$(iii)$] $U^*UW_{i}=S_{\mathcal{W}}^{-1}W_{i}$ for all $ i \in I.$   
\item[$(iv)$] $\{(UW_{i},\omega_{i})\}_{i \in I} $ is an orthogonal fusion Riesz basis.
\end{itemize}
\end{theorem}

\begin{proof}
$(i)\Leftrightarrow (iii)$ Let $\mathcal{W}$ be $\left( U,\omega^{-1}\right) $-scalable. Using the fact that
\begin{equation}\label{tyt}
UW_{i}\perp \left( U^*\right) ^{-1}S_{\mathcal{W}}^{-1}W_{j},\quad (i\neq j),
\end{equation}
we get
\begin{equation}\label{sigma2}
\begin{aligned}
\left( U^*\right) ^{-1}S_{\mathcal{W}}^{-1}\pi_{W_{i}}&=S_{U\mathcal{W}_{\omega^{-1}}}\left( U^*\right) ^{-1}S_{\mathcal{W}}^{-1}\pi_{W_{i}}
\\&=\sum_{j \in I}\pi_{UW_{j}}\left( U^*\right) ^{-1}S_{\mathcal{W}}^{-1}\pi_{W_{i}}
\\&=\pi_{UW_{i}}\left( U^*\right) ^{-1}S_{\mathcal{W}}^{-1}\pi_{W_{i}}, \quad (i \in I).
\end{aligned}
\end{equation}
Hence, $ \left( U^*\right) ^{-1}S_{\mathcal{W}}^{-1}W_{i}\subseteq UW_{i} ,$ which implies that $S_{\mathcal{W}}^{-1}W_{i}\subseteq U^*UW_{i},$ for all $i \in I$.
On the other hand, by applying \eqref{sigma2} and Lemma \ref{r.e}$(iv)$, we get
\begin{equation*}
\left( U^*\right) ^{-1}S_{\mathcal{W}}^{-1}\pi_{W_{i}}=\pi_{UW_{i}}\left( U^*\right) ^{-1}S_{\mathcal{W}}^{-1}\pi_{W_{i}}=\pi_{UW_{i}}\left( U^*\right) ^{-1}.
\end{equation*}
So, for each $f,g \in \mathcal {H},$ we induce
\begin{align*}
\left\langle  \pi_{UW_{i}}US_{\mathcal{W}}f,g \right\rangle  &= \left\langle f , S_{\mathcal{W}}U^{*}\pi_{UW_{i}}g \right\rangle 
\\&= \left\langle f , S_{\mathcal{W}}U^{*}\pi_{UW_{i}}\left( U^*\right) ^{-1}U^*g \right\rangle 
\\&= \left\langle f , \pi_{W_{i}}U^*g \right\rangle  = \left\langle  U\pi_{W_{i}}f,g \right\rangle.
\end{align*}
Thus, $\pi_{UW_{i}}US_{\mathcal{W}}=U\pi_{W_{i}},$ for all $i \in I.$ It follows from Lemma \ref{extend Gav} that $U^*UW_{i}\subseteq S_{\mathcal{W}}^{-1}W_{i}.$ Therefore, $U^*UW_{i}=S_{\mathcal{W}}^{-1}W_{i},$ for all $i \in I$. For the converse implication, since $UW_{i}=\left( U^*\right) ^{-1}S_{\mathcal{W}}^{-1}W_{i},$ then by employing Lemma \ref{r.e}, \eqref{tyt} and taking $ \gamma_{i}=\omega_{i}^{-1},$ we conclude that
\begin{align*}
S_{U\mathcal{W}_{\gamma}}\left( U^*\right) ^{-1}&=\sum_{i \in I}\pi_{UW_{i}}\left( U^*\right) ^{-1}
\\&=\sum_{i \in I}\omega_{i}^{2}\pi_{UW_{i}}\left( U^*\right) ^{-1}S_{\mathcal{W}}^{-1}\pi_{W_{i}}
\\&=\sum_{i \in I}\omega_{i}^{2}\pi_{\left( U^*\right) ^{-1}S_{\mathcal{W}}^{-1}W_{i}}\left( U^*\right) ^{-1}S_{\mathcal{W}}^{-1}\pi_{W_{i}}
\\&=\left( U^*\right) ^{-1}\sum_{i \in I}\omega_{i}^{2}S_{\mathcal{W}}^{-1}\pi_{W_{i}}=\left( U^*\right) ^{-1}.
\end{align*}
Hence, $S_{U\mathcal{W}_{\gamma}}=I_{\mathcal {H}}$ and consequently $\mathcal{W}$ is $\left( U,\omega^{-1}\right) $-scalable.

$(i)\Leftrightarrow (ii)$ Let $\mathcal{W}$ be $\left( U,\omega^{-1}\right) $-scalable. In view of $(iii)$, we get
\begin{align*}
\left( US_{\mathcal{W}}\right) ^{-1}\left( S_{\mathcal{W}}U^*\right) ^{-1}W_{i}&=S_{\mathcal{W}}^{-1}U^{-1}\left( U^*\right) ^{-1}S_{\mathcal{W}}^{-1}W_{i}
\\&=S_{\mathcal{W}}^{-1}U^{-1}UW_{i}
\\&=S_{\mathcal{W}}^{-1}W_{i},
\end{align*}
for each $ i\in I $. Therefore, $\mathcal{W}$ is $\left( S_{\mathcal{W}}U^*\right) ^{-1}$-scalable by using $(iii)\Rightarrow(i)$. The converse is derived by applying a symmetry argument.



$(iii)\Rightarrow (iv)$ Applying $(iii)$ leads to 
\begin{equation*}
\left\langle UW_{i},UW_{j} \right\rangle = \left\langle U^{*}UW_{i},W_{j} \right\rangle = \left\langle S_{\mathcal{W}}^{-1}W_{i},W_{j} \right\rangle =0,
\end{equation*}
for each $i\neq j.$ Thus, $(iv)$ holds.


$(iv)\Rightarrow (i)$ Taking $ \gamma_{i}=\omega_{i}^{-1}$, we obtain
\begin{equation*}
S_{U\mathcal{W}_{\gamma}}=\sum_{i \in I} \pi_{UW_{i}}=\pi_{\oplus_{i \in I} UW_{i}}=I_{\mathcal {H}},
\end{equation*}
which yields $(i)$ is satisfied.
\end{proof}

\end{document}
